\chapter{On a global correspondence}

\section{Introduction}
In this chapter, we use results from previous chapters
to obtain a correspondence using the trace
formula between cuspidal automorphic representations of
$G({\Bbb {A}})=SL(2,{\Bbb {A}})$ and genuine cuspidal
automorphic representations of the 2-fold
metaplectic cover $\overline{SL(2,{\Bbb {A}})}$.

Our main result here is the following

\begin{theorem}
There is a correspondence between automorphic cuspidal representations
of $G({\Bbb A})$ and those of $\overline{G({\Bbb A})}$:

(a) to each cuspidal automorphic representation $\pi $ of $G({\Bbb A})$
there is a {\em unique} genuine cuspidal automorphic
representation $\overline{\pi }$ of
$\overline{G({\Bbb A})}$ such that 
\[
{\rm tr}\;\pi _v(f_v)={\rm tr}\;\overline{\pi }_v(\overline{f}_v)
\]
for almost all places $v$, where $f_v$ and $\overline{f}_v$ are
corresponding in the sense described below,

(b) to each cuspidal automorphic representation
$\overline{\pi }$ of $\overline{G({\Bbb A})}$ there
is a genuine cuspidal automorphic
representation $\pi $ of $G({\Bbb A})$ such that 
\[
{\rm tr}\;\pi _v(f_v)={\rm tr}\;\overline{\pi }_v(\overline{f}_v)
\]
for almost all places $v$, where $f_v$ and $\overline{f}_v$ are
corresponding in the sense described below.

Moreover, a cuspidal automorphic representation $\pi $ of $G({\Bbb A})$
occurs with multiplicity one provided

(i) $\pi $ corresponds to a $\overline{\pi }$ as in (a), and

(ii) ${\rm tr}\;\pi _v(f_v)={\rm tr}\;\overline{\pi }_v(\overline{f}_v)$,
for {\em all} places $v$, where $f_v$ and $\overline{f}_v$
are corresponding in the sense described below.
\end{theorem}

\begin{definition} Let $p$ be a place of ${\Bbb{Q}}$.
We say that an automorphic representation
$\pi$ is {\bf $p$-good}
if $\pi=\otimes \pi_v$ and $p|v$ then $\pi_v$ is
either a unitary principal series (in the ``continuous series'') or
a complementary series.
\end{definition}

\begin{corollary} If $\pi=\otimes_v \pi_v$ is a 2-good
automorphic cuspidal representation of $G({\Bbb{A}})$ then $\pi$
occurs with multiplicity one in the spectrum of
$L^2(G({\Bbb{Q}})\backslash G({\Bbb{A}}))$.
\end{corollary}

proof: This corollary is a consequence of the proof below and the local
comparisons in previous chapters but some words of
explanation are warranted.
The result is a straightforward consequence of
the last part of the theorem and the character formulas
of chapters 4, 5, 6, 7 for all automorphic representations
which are either $p$-good for all $p$ or have a supercuspidal
component at some place $v$ not dividing $2$. (It
may very well follow from \cite{LL} in this case as well.)

We first focus our attention on the $\pi=\otimes_w\pi_w$
with a local component $\pi_v$, $v$ not dividing $2$, which is
special: $\pi_v=\sigma_v$ (the special representation,
for $v$ non-archimedian, and the ``lowest'' discrete series
whose character is equal to $-1$ on $SO(2,{\Bbb{R}})$, for
$v$ real). We know from chapter 6, \S 5, that
if $v$ is non-archimedian then there is a matching
$tr \sigma_v(f_v)=tr\overline{\sigma}_v(\overline{f}_v)$,
where $f_v$, $\overline{f}_v$ are matching and are
supported in a neighborhood of the identity. Here
$\overline{\sigma}_v$ denotes the odd oscillator representation
$\omega_{odd}$. The archimedian case follows from
the character relations in \cite{Fl}, \S 2.4, for functions
supported on the elliptic set. See J. Adams \cite{Ad}
for the character of $\overline{\sigma}_v$ in general
for $v$ real. Similarly for the trivial representation.

We must extract a result, proven in the next section,
which follows from the comparison of the
trace formulas. Fix a place $v_0$.
Let $V$ be a finite set of places
containing at most one place $w$ where $\pi_w$ is a discrete series
and containing $v_0$ but no other place $w$ where $\pi_w$ is
a reducible principal series. Then
\[
\sum\Sb \pi^{\prime }\\ \pi _w^{\prime }\cong\pi _w^0\\ w\notin V \endSb 
m(\pi^{\prime })\prod_{v\in V}{\rm tr}\;\pi _v^{\prime }(f_v)
=\prod_{v\in V}{\rm tr}\;\overline{\pi }_v(\overline{f}_v).
\]
If $V=\{v_0\}$ then the above equation implies
\begin{equation}
\sum\Sb \pi^{\prime }\\ \pi _w^{\prime }\cong\pi _w^0\\ w\not= v_0 \endSb 
m(\pi^{\prime }){\rm tr}\;\pi _{v_0}^{\prime }(f_{v_0})
={\rm tr}\;\overline{\pi }_{v_0}(\overline{f}_{v_0}).
\end{equation}

Let $v_0$ be real.
If $\overline{\pi}_{v_0}=\overline{\sigma}_{v_0}$ and
if $\overline{f}_{v_0}$, $f_{v_0}$ are matching and supported on the
elliptic set then this last equation and the above remarks imply
that there are constants $a_k$ such that
\[
-\overline{F}^{\wedge}(0)=\sum_k a_k F^{\wedge}(k),
\]
where $F(\theta)=F_{f_{v_0}}(r(\theta))$,
where $r(\theta)\in SO(2,{\Bbb{R}})$
is the usual rotation matrix (as in \cite{Fl}, \S 2.4, for
example), where
$\overline{F}(\theta)=F_{\overline{f}_{v_0}}(r(\theta))$, and where
$F^\wedge$ denote the Fourier transform with respect to $\theta$.
Since $f_{v_0}$, $\overline{f}_{v_0}$ are matching,
$\overline{F}^\wedge(0)=F^\wedge(0)$. By linear independence, this forces
$a_0=-1$ and $a_k=0$ for $k>1$. This implies that 
$tr \sigma_v(f_{v_0})=tr\overline{\sigma}_{v_0}(\overline{f}_{v_0})$,
for all $f_v$, $\overline{f}_v$ which are matching. The result now
follows for $\pi$ with $\sigma_v$ as a local constituent.

Since the character formulas in chapters 5, 7
imply that the characters of irreducible
admissible representations are linearly independent in a
neighborhood of the identity, the last part of the theorem holds
if condition (ii) is replaced by a local matching condition.
This, in conjunction with the proof of the theorem below,
implies the corollary for all automorphic representations
which do not contain local components $\pi_v$ which are
limits of discrete series. 

The limits of discrete series representations do not have a
matching as in part (ii) of the theorem because there
are no reducible unitary principal series on the cover to match
with (at least if $p\not= 2$, see \S 6.3.1).
Let $\pi=\otimes_v\pi_v$ be a cuspidal automorphic
representation and let $\pi_{v_0},\pi_{v_0}'$ be as in the statement
of the corollary below (we shall finish the proof of the above
corollary and prove the corollary below at the same time).

We know that $tr\overline{\pi}_{v_0}(\overline{f}_v)
=tr \pi_{v_0}(f_{v_0})+tr \pi'_{v_0}(f_{v_0})$,
where $\overline{\pi_{v_0}}$ is the representation
corresponding (as in chapters 4,5,6)
to the full induced unitary representation containing $\pi_{v_0}$
as a component.
This fact, along with linear independence and the equation
(10.1), implies Corollaries 345 above and 346 below. $\Box$

\begin{corollary} Let $\pi=\otimes_v \pi_v$ be an automorphic
cuspidal representation of $G({\Bbb{A}})$.
Let $w$ be a non-archimedian place of $k$ with
$|2|_w=1$, let $\pi_w$ be a reducible unitary principal
series (a ``limit of
discrete series'') of $SL(2,k_w)$ occuring as a local component
of $\pi$ and let $\pi'_w$ is its
``complement'' (in the full induced unitary representation).
Then $\pi'=\otimes_{v\not= w}\pi_v\otimes\pi'_w$
is a cuspidal automorphic representation.
\end{corollary}

The previous corollary is very similar to results in
\cite{LL} but I don't see how to derive it from their
results.

\begin{definition} Let $V$ be a finite set of places and let
$\pi =\otimes _v\pi _v,\ \pi' =\otimes _v\pi' _v$ be two
representations of $G({\Bbb{A}})$.
We say that $\pi$ is {\bf $V$-nearly equivalent} to $\pi'$ if
$\pi_v\cong \pi'_v$ for all $v\notin V$.
\end{definition}

We shall make the correspondence more specific.

\begin{corollary}
Let $\pi =\otimes _v\pi _v$ be a fixed cuspidal automorphic representation
of $G({\Bbb A})$.
Let $V=V(\pi)$ be a finite set of places containing at most one
place $v$ for which $\pi_v$ is a discrete series.
For each $w\notin V$, let $\pi _w^0=\pi _w$. There is a
unique genuine cuspidal automorphic representation $\overline{\pi }=\otimes
_v\overline{\pi }_v$ of $\overline{G({\Bbb A})}$ such that 
\[
\sum\Sb \pi ^{\prime } \\ \pi _w^{\prime }\cong \pi _w^0 \\
w\notin V \endSb
m(\pi ^{\prime })\prod_{v\in V}
{\rm tr}\;\pi _v^{\prime }(f_v)
=\prod_{v\in V}{\rm tr}\;\overline{\pi }_v(\overline{f}_v),
\]
where one of the representations $\pi ^{\prime }$ occurring in the sum on
the left is $\pi $.

Conversely, given a genuine cuspidal automorphic
representation $\overline{\pi }=\otimes _v\overline{\pi }_v$ of
$\overline{G({\Bbb A})}$, it corresponds to a cuspidal automorphic
representation $\pi $
of $G({\Bbb A})$ if and only if $\pi $ is $V$-nearly equivalent to
some representation $\pi^0=\otimes _v \pi^0_v$ of
$G({\Bbb{A}})$, for some set of places $V$ with the following
two properties:

(i) the complement $V'$ of $V$ contains all but possibly one of the
the places $w$ for which $\overline{\pi}_w$ is a discrete
series representation, and

(ii) $\overline{\pi}_v$ corresponds to $\pi_v^0$ for all $v\notin V$.
\end{corollary}

\section{Proof of the main result}

We write the trace formula of $\overline{G({\Bbb A})}$ in
the form
\begin{eqnarray}
&&{\rm tr}\;R_{cusp}(\phi )-S_1(\phi )+
\sum_{\chi \in \overline{{\cal X}}-\overline{{\cal X}}(G)}
I_\chi (\phi ) \\
&=&\sum_{{\cal O}\in {\cal C}}SI_{{\cal O}}(\phi )
+S_2(\phi )+S_3(\phi )+
\sum\Sb \Sb H \\ cuspidal \endSb  \\ endoscopic \endSb
\sum_{{\cal O}_H}SI_{{\cal O}_H}^H(\phi ^H).  \nonumber
\end{eqnarray}
We write the trace formula of $G({\Bbb A})$ in the form
\begin{eqnarray}
&&\ {\rm tr}\;R_{cusp}(f)-S_1(f)+
\sum_{\chi \in {\cal X}-{\cal X}(G)}I_\chi(f) \\
\  &=&\sum_{{\cal O}\in {\cal C}}SI_{{\cal O}}(f)
+S_2(f)+S_3(f)+\sum\Sb \Sb H
\\ cuspidal \endSb  \\ endoscopic \endSb
\sum_{{\cal O}_H}SI_{{\cal O}_H}^H(f^H).  \nonumber
\end{eqnarray}

This trace formula differs slightly from that in \cite{LL}
in that its ``spectral
side'' has not been put in stable form. This form is more
convenient for us.

Assume that $f,\phi $ are both factorizable. We shall replace $\phi $ by a
genuine function $\overline{f}$ such that
$\overline{f}_v$ matches with $f_v$ for all $v$ in the sense of
chapter 5, \S 3. Since we have only matched weighted orbital integrals
of spherical functions at the unit elements, we must assume that
at each place $f$ and $\overline{f}$ are either both
smooth functions or both unit elements. Moreover, at the real place
and at the places dividing $2$, we assume that they are smooth.

It turns out to be relatively simple (see chapter 5)
to compare all of the
corresponding terms between these two expressions, except for
\begin{equation}
{\rm tr}\;R_{cusp}(\overline{f})\stackrel{?}{=}{\rm tr}\;R_{cusp}(f)
\end{equation}
and
\begin{equation}
S_1(\overline{f})\stackrel{?}{=}S_1(f)
\end{equation}
and 

\begin{equation}
\sum_{\chi \in \overline{{\cal X}}-\overline{{\cal X}}(G)}
I_\chi ^{\prime }(\overline{f})\stackrel{?}{=}
\sum_{\chi \in {\cal X}-{\cal X}(G)}I_\chi^{\prime }(f),
\end{equation}
and the stable $A$-elliptic terms in \S 8.6.9 (c), \S 9.5.8 (c).
\begin{remark}
All we really care about is proving the first equation (10.4). One idea,
originally due to Langlands \cite{L4}, is to prove that 
\[
{\rm tr}\;R_{cusp}(\overline{f})-{\rm tr}\;R_{cusp}(f)=S_1(\overline{f}%
)-S_1(f)-\sum_{\chi \in \overline{{\cal X}}-\overline{{\cal X}}(G)}I_\chi
^{\prime }(\overline{f})+\sum_{\chi \in {\cal X}-{\cal X}(G)}I_\chi ^{\prime
}(f)
\]
cannot be true (which we know it is by the comparisons mentioned above)
unless both sides are zero. It involves facts about the Satake transform and
was used in \cite{L4}, \cite{Fl}, sections 4.1-2
and \cite{R}, sections 10.5, 14.5 (where it is called the
decomposition of measures argument).

Another idea is that used in \cite{AC}, ch 2, section 10.
It may apply but it is
so technical I haven't been able to understand it fully yet.

We shall use another idea. It is described in chapter 5
and relies on results in \cite{A6}.
\end{remark}

The comparison in chapter 5,
immediately implies the matching for $S_1$ in (10.5) is true.

Recall that
\begin{equation}
I_\chi ^{\prime }(\overline{f})=
\sum_{\overline{\eta }\in \overline{E}_\chi }
\frac{m^{\prime }(\overline{\eta })}{m(\overline{\eta })}
{\rm tr}\;I_A(\overline{\eta },\overline{f})
\end{equation}
and 
\begin{equation}
I_\chi ^{\prime }(f)=\sum_{\eta \in E_\chi }\frac{m^{\prime
}(\eta )}{m(\eta )}{\rm tr}\;I_A(\eta ,f),
\end{equation}
where the indexing sets $\overline{E}_\chi ,E_\chi $ have been described
previously. 

We claim that the indexing sets $\overline{E}_\chi ,E_\chi $ are in 1-1
correspondence.

We also claim that the indexing sets
$\overline{{\cal X}}-\overline{{\cal X}}(G)$,
${\cal X}-{\cal X}(G)$ are in 1-1 correspondence.

Both of these claims follow from that fact that
$\overline{A({\Bbb{A}}_S)}$ is
abelian.

The comparison between ${\rm tr}\;I_A(\overline{\eta },\overline{f})$,
${\rm tr}\;I_A(\eta ,f)$ was done in chapter 5. 

{\bf Claim}
If $\overline{\chi }\in \overline{{\cal X}}-\overline{{\cal X}}(G)$
corresponds to $\chi \in {\cal X}-{\cal X}(G)$ and if $\overline{\eta }\in 
\overline{E}_{\overline{\chi }}$ corresponds to $\eta \in E_\chi $ then
$m(\overline{\eta })=m(\eta )$.

This follows along the same lines as the analogous comparison
in \cite{Fl}. We use chapter 5, at the finite places and
\cite{Sa} at the real places. At the complex archimedian places (if any),
the local covering splits and known results apply
(see \cite{Sha} for example).

To compare (10.7), (10.8), and the stable hyperbolic terms,
it suffices to compare
\[
I_A(\pi ,f)\;\;{\rm and}\;\;I_A(\overline{\pi },\overline{f})\;
\]
and
\[
\varphi _A^G(\pi ,f)\;\;{\rm and}\;\;
\varphi _A^G(\overline{\pi },\overline{f}),\;
\]
for $(\pi ,A)\in {\cal X}-{\cal X}(G)$ corresponding to
$(\overline{\pi },A)\in \overline{{\cal X}}-\overline{{\cal X}}(G)$.
These boil down to comparisons established in chapter 5. 

We conclude that (10.4) is true, as desired.

We now proceed to use (10.4) to verify our main result.

It is known from Waldspurger \cite{Wa1} (see also
\cite{GPS} and \cite{PS})
that $m(\overline{\pi })=0$ or $=1$ for all genuine
automorphic cuspidal representations $\overline{\pi }$ of
$\overline{G({\Bbb A})}$. However, strong
multiplicity one fails to hold for $\overline{G({\Bbb A})}$.
There is a weak (pun intended) version of strong
multiplicity one for $\overline{G({\Bbb A})}$ due to Waldspurger
whose form below I derive from \cite{PS}:

\begin{lemma} (\cite{Wa2}) Let $\overline{\pi }
=\otimes_v\overline{\pi }_v$ 
be a genuine automorphic cuspidal representations of
$\overline{G({\Bbb A})}$. Let $V$ be a finite set of places. Let
\[
%\begin{array}{c}
\Sigma=\Sigma(\overline{\pi })=\{v\ |\ 
\overline{\pi }_v\ {\rm is \ a \ discrete\ series\ representation}
%\ \ \ \ \ \ \ \ \ \ \ \ \ \ \ \
%{\rm arising \ from\ the\ correspondence\ from}\ O(1)
\}.
%\end{array}
\]
Assume $\Sigma\cap V$ has even cardinality.
The number of genuine automorphic cuspidal representations
$\overline{\pi}'$ of $\overline{G({\Bbb A})}$ such that
\[
\overline{\pi }_v\cong \overline{\pi }_v',\ \ \ \
{\rm for\ all\ }v\notin V,
\]
is $1$ if $\Sigma \subset V$ and $2^{|\Sigma|-|\Sigma\cap V|-1} $
otherwise.
\end{lemma}

Let $V$ denote a finite set of places of $k$ containing all
$v\in \Sigma^\infty_k$.

Let $w\notin V$ be fixed for the moment and let $\pi _w^0$,
$\overline{\pi }_w^0$ be fixed corresponding
unramified principal series representations (in
the sense of chapter 5).
We claim that
\[
\sum\Sb \pi  \\ \pi _w\cong \pi _w^0 \endSb m(\pi )
\prod_{v\neq w}{\rm tr}\;
\pi _v(f_v)=\sum\Sb \overline{\pi } \\
\overline{\pi }_w\cong \overline{\pi }_w^0 \endSb
m(\overline{\pi })\prod_{v\neq w}{\rm tr}\;
\overline{\pi }_v(\overline{f}_v).
\]
This is the analog of \cite{GJ}, Lemma 8.18 and is proven the same way.

Now, for each $w\notin V$, fix corresponding representations
$\pi _w^0$, $\overline{\pi }_w^0$ in the sense of chapter 5 or chapter 7,
almost all of which
must be unramified principal series representations.
We claim that
\[
\sum\Sb \pi  \\ \pi _w\cong \pi _w^0 \\ w\notin V \endSb m(\pi )\prod_{v\in
V}{\rm tr}\;\pi _v(f_v)=\sum\Sb \overline{\pi } \\ \overline{\pi }_w\cong 
\overline{\pi }_w^0 \\ w\notin V \endSb
m(\overline{\pi })\prod_{v\in V}{\rm tr}\;
\overline{\pi }_v(\overline{f}_v).
\]
This is the analog of \cite{GJ}, Lemma 8.19 and is proven the same way.

Let $\pi =\otimes _v\pi _v$ be a fixed cuspidal automorphic representation
of $G({\Bbb A})$. For each $w\notin V$, let $\pi _w^0=\pi _w$.
Assume $V$ has been choosen so that its complement
contains all, with possibly one exception, places $w$ such that
$\pi_w^0$ is a discrete series. Each $\pi_w^0$, $w\notin V$,
which is a discrete series corresponds to an
$\overline{\pi}_w^0$ which is also a discrete series.
We claim that
there is a unique genuine cuspidal automorphic representation $\overline{\pi 
}=\otimes _v\overline{\pi }_v$ of $\overline{G({\Bbb A})}$ such that 
\begin{equation}
\sum\Sb \pi ^{\prime } \\ \pi _w^{\prime }\cong \pi _w^0 \\ w\notin V \endSb %
m(\pi ^{\prime })\prod_{v\in V}{\rm tr}\;\pi _v^{\prime }(f_v)=\prod_{v\in V}%
{\rm tr}\;\overline{\pi }_v(\overline{f}_v).
\end{equation}
This is the analog of \cite{GJ}, Proposition 8.20
and is proven the same way. (We have used Lemma 349
above to obtain uniqueness and multiplicity one to
take $m(\overline{\pi })=1$.) This proves part (a) of the theorem above.

Part (b) of the theorem is proven similarly. (Note in this case, neither
strong multiplicity one or multiplicity one may be invoked.)

To verify the last part, by our hypothesis, with $V$ choosen as above,
\[
\sum\Sb \pi ^{\prime } \\ \pi _w^{\prime }\cong \pi _w^0 \\
w\notin V \endSb
m(\pi ^{\prime })\prod_{v\in V}{\rm tr}\;
\pi _v^{\prime }(f_v)
=\prod_{v\in V}{\rm tr}\;\overline{\pi }_v(\overline{f}_v)
=\prod_{v\in V}{\rm tr}\;\pi_v(f_v),
\]
so 
\[
\sum\Sb \pi ^{\prime }\neq \pi  \\ \pi _w^{\prime }\cong \pi _w^0 \\ w\notin
V \endSb m(\pi ^{\prime })\prod_{v\in V}{\rm tr}\;\pi _v^{\prime
}(f_v)+(m(\pi )-1)\prod_{v\in V}{\rm tr}\;\pi _v(f_v)=0.
\]
By the linear independence of characters, $m(\pi ^{\prime })=0$ for $\pi
^{\prime }\neq \pi $ and $m(\pi )-1=0$, as desired. $\Box $


